%%%%%%%%%%%%%%%%%%%%%%%%%%%%%%%%%%%%%%%%%%%%%%%%%%%%%%%%%%%%
% 7.13 STOCHASTIC CALCULUS
% The Composition of Advanced Mathematics
%%%%%%%%%%%%%%%%%%%%%%%%%%%%%%%%%%%%%%%%%%%%%%%%%%%%%%%%%%%%

\section{Stochastic Calculus}
\label{sec:stochastic-calculus}

\prerequisite{
Measure-theoretic probability, stochastic processes, Brownian motion,
filtrations, conditional expectation, martingales, Lebesgue integration,
multivariable calculus, Taylor expansions, and ordinary differential
equations.
}

\learningobjectives{
By the end of this section, the reader should be able to:
\begin{itemize}
    \item explain why ordinary calculus is insufficient for Brownian
          sample paths;
    \item define quadratic variation;
    \item compute the quadratic variation of Brownian motion;
    \item interpret the heuristic rule
          \[
          (dB_t)^2=dt;
          \]
    \item define simple adapted stochastic processes;
    \item define the It\^o stochastic integral for simple processes;
    \item understand construction of the It\^o integral by \(L^2\)
          approximation;
    \item apply the It\^o isometry;
    \item compute expectations and variances of stochastic integrals;
    \item recognize It\^o integrals as martingales;
    \item distinguish ordinary differentials from stochastic
          differentials;
    \item state and apply It\^o's formula;
    \item understand the multidimensional It\^o formula;
    \item derive common stochastic differentials;
    \item define quadratic covariation;
    \item use stochastic product and integration-by-parts formulas;
    \item understand exponential martingales;
    \item distinguish It\^o and Stratonovich integrals;
    \item convert between It\^o and Stratonovich notation in elementary
          settings;
    \item understand localization and local martingales conceptually;
    \item recognize stochastic exponentials;
    \item understand change of measure and Girsanov's theorem at an
          introductory level;
    \item connect stochastic calculus with stochastic differential
          equations, finance, diffusion, filtering, and physics.
\end{itemize}
}

%%%%%%%%%%%%%%%%%%%%%%%%%%%%%%%%%%%%%%%%%%%%%%%%%%%%%%%%%%%%
\subsection{Why Ordinary Calculus Is Not Enough}
\label{subsec:why-ordinary-calculus-not-enough}
%%%%%%%%%%%%%%%%%%%%%%%%%%%%%%%%%%%%%%%%%%%%%%%%%%%%%%%%%%%%

Ordinary differential calculus is designed for sufficiently smooth
functions.

Brownian motion does not have smooth sample paths.

A standard Brownian path is almost surely:

\[
\boxed{
\text{continuous}
}
\]

but also

\[
\boxed{
\text{nowhere differentiable}.
}
\]

Therefore an expression such as

\[
\frac{dB_t}{dt}
\]

does not exist in the ordinary pathwise sense.

Nevertheless, one would like to make mathematical sense of expressions
such as

\[
\int_0^t H_s\,dB_s.
\]

Stochastic calculus provides the framework for doing so.

%%%%%%%%%%%%%%%%%%%%%%%%%%%%%%%%%%%%%%%%%%%%%%%%%%%%%%%%%%%%
\subsection{Brownian Motion Review}
\label{subsec:brownian-review-stochastic-calculus}
%%%%%%%%%%%%%%%%%%%%%%%%%%%%%%%%%%%%%%%%%%%%%%%%%%%%%%%%%%%%

A standard Brownian motion

\[
(B_t)_{t\geq0}
\]

satisfies:

\[
\boxed{
B_0=0,
}
\]

has independent increments,

has stationary increments,

has continuous sample paths almost surely, and

\[
\boxed{
B_t-B_s
\sim
N(0,t-s)
}
\]

for

\[
0\leq s<t.
\]

Its covariance is

\[
\boxed{
\operatorname{Cov}(B_s,B_t)
=
\min(s,t).
}
\]

%%%%%%%%%%%%%%%%%%%%%%%%%%%%%%%%%%%%%%%%%%%%%%%%%%%%%%%%%%%%
\subsection{Size of Brownian Increments}
\label{subsec:size-brownian-increments}
%%%%%%%%%%%%%%%%%%%%%%%%%%%%%%%%%%%%%%%%%%%%%%%%%%%%%%%%%%%%

For a small time increment

\[
\Delta t,
\]

Brownian motion satisfies

\[
\Delta B
\sim
N(0,\Delta t).
\]

Therefore its typical magnitude is of order

\[
\boxed{
\sqrt{\Delta t}.
}
\]

This differs fundamentally from a differentiable function, whose
increment is typically of order

\[
\Delta t.
\]

Thus Brownian motion fluctuates much more strongly on very small time
scales than a smooth deterministic path.

%%%%%%%%%%%%%%%%%%%%%%%%%%%%%%%%%%%%%%%%%%%%%%%%%%%%%%%%%%%%
\subsection{Quadratic Variation}
\label{subsec:quadratic-variation}
%%%%%%%%%%%%%%%%%%%%%%%%%%%%%%%%%%%%%%%%%%%%%%%%%%%%%%%%%%%%

Consider a partition

\[
0=t_0<t_1<\cdots<t_n=t.
\]

For a function or process \(X\), define the quadratic-variation sum

\[
\boxed{
\sum_{k=0}^{n-1}
\left(
X_{t_{k+1}}-X_{t_k}
\right)^2.
}
\]

As the partition becomes finer, this sum may converge.

\begin{definitionbox}{Quadratic Variation}
The quadratic variation of \(X\) on \([0,t]\), when the limit exists,
is denoted

\[
\boxed{
[X]_t
}
\]

and is the appropriate limit of

\[
\sum_k
(\Delta X_k)^2
\]

as the partition mesh tends to zero.
\end{definitionbox}

%%%%%%%%%%%%%%%%%%%%%%%%%%%%%%%%%%%%%%%%%%%%%%%%%%%%%%%%%%%%
\subsection{Quadratic Variation of Smooth Functions}
\label{subsec:qv-smooth-functions}
%%%%%%%%%%%%%%%%%%%%%%%%%%%%%%%%%%%%%%%%%%%%%%%%%%%%%%%%%%%%

Suppose \(f\) is continuously differentiable.

Then

\[
\Delta f
\approx
f'(t)\Delta t.
\]

Therefore,

\[
(\Delta f)^2
\]

is of order

\[
(\Delta t)^2.
\]

Summing approximately \(1/\Delta t\) such terms gives a total of order

\[
\Delta t.
\]

Hence,

\[
\boxed{
[f]_t=0
}
\]

for sufficiently smooth deterministic functions.

%%%%%%%%%%%%%%%%%%%%%%%%%%%%%%%%%%%%%%%%%%%%%%%%%%%%%%%%%%%%
\subsection{Quadratic Variation of Brownian Motion}
\label{subsec:qv-brownian}
%%%%%%%%%%%%%%%%%%%%%%%%%%%%%%%%%%%%%%%%%%%%%%%%%%%%%%%%%%%%

Brownian increments satisfy

\[
\Delta B
\sim
\sqrt{\Delta t}.
\]

Therefore,

\[
(\Delta B)^2
\]

is of order

\[
\Delta t.
\]

Summing over approximately \(1/\Delta t\) intervals gives a quantity of
order \(1\).

In fact,

\[
\boxed{
[B]_t=t
}
\]

almost surely in the standard quadratic-variation sense.

This is one of the central facts of stochastic calculus.

%%%%%%%%%%%%%%%%%%%%%%%%%%%%%%%%%%%%%%%%%%%%%%%%%%%%%%%%%%%%
\subsection{Worked Example: Expected Quadratic Variation}
\label{subsec:worked-expected-qv}
%%%%%%%%%%%%%%%%%%%%%%%%%%%%%%%%%%%%%%%%%%%%%%%%%%%%%%%%%%%%

\begin{workedexamplebox}{Brownian Partition Sum}

Take an equally spaced partition of \([0,t]\) into \(n\) intervals.

Then

\[
\Delta t
=
\frac{t}{n}.
\]

Each Brownian increment satisfies

\[
\mathbb{E}
[
(\Delta B_k)^2
]
=
\frac{t}{n}.
\]

Therefore,

\[
\mathbb{E}
\left[
\sum_{k=0}^{n-1}
(\Delta B_k)^2
\right]
=
n\frac{t}{n}.
\]

\begin{answerbox}
\[
\boxed{
\mathbb{E}
\left[
\sum_{k=0}^{n-1}
(\Delta B_k)^2
\right]
=
t.
}
\]
\end{answerbox}

This is consistent with

\[
[B]_t=t.
\]

\end{workedexamplebox}

%%%%%%%%%%%%%%%%%%%%%%%%%%%%%%%%%%%%%%%%%%%%%%%%%%%%%%%%%%%%
\subsection{The Heuristic Differential Rules}
\label{subsec:ito-differential-rules}
%%%%%%%%%%%%%%%%%%%%%%%%%%%%%%%%%%%%%%%%%%%%%%%%%%%%%%%%%%%%

Quadratic variation motivates the symbolic rules

\[
\boxed{
(dB_t)^2=dt,
}
\]

\[
\boxed{
dB_t\,dt=0,
}
\]

and

\[
\boxed{
(dt)^2=0.
}
\]

These are not ordinary algebraic identities.

They are shorthand for the scaling and quadratic-variation structure
that appears when stochastic differentials are manipulated through
It\^o's formula.

%%%%%%%%%%%%%%%%%%%%%%%%%%%%%%%%%%%%%%%%%%%%%%%%%%%%%%%%%%%%
\subsection{Why \((dB_t)^2\) Survives}
\label{subsec:why-db-square-survives}
%%%%%%%%%%%%%%%%%%%%%%%%%%%%%%%%%%%%%%%%%%%%%%%%%%%%%%%%%%%%

Since

\[
dB_t
\]

behaves on the scale

\[
\sqrt{dt},
\]

its square behaves on the scale

\[
dt.
\]

Thus second-order Brownian terms contribute at the same order as
first-order time terms.

This is the crucial difference from ordinary calculus.

It is why stochastic Taylor expansions retain a second derivative term.

%%%%%%%%%%%%%%%%%%%%%%%%%%%%%%%%%%%%%%%%%%%%%%%%%%%%%%%%%%%%
\subsection{Simple Adapted Processes}
\label{subsec:simple-adapted-processes}
%%%%%%%%%%%%%%%%%%%%%%%%%%%%%%%%%%%%%%%%%%%%%%%%%%%%%%%%%%%%

Let

\[
0=t_0<t_1<\cdots<t_n=T.
\]

A simple adapted process has the form

\[
\boxed{
H_t
=
\sum_{k=0}^{n-1}
H_k
\mathbf{1}_{(t_k,t_{k+1}]}(t),
}
\]

where each

\[
H_k
\]

is measurable with respect to

\[
\mathcal{F}_{t_k}.
\]

Thus the value used on an interval is determined by information
available at the beginning of that interval.

%%%%%%%%%%%%%%%%%%%%%%%%%%%%%%%%%%%%%%%%%%%%%%%%%%%%%%%%%%%%
\subsection{Why Adaptedness Matters}
\label{subsec:why-adaptedness-ito}
%%%%%%%%%%%%%%%%%%%%%%%%%%%%%%%%%%%%%%%%%%%%%%%%%%%%%%%%%%%%

The stochastic integral is intended to model quantities chosen without
knowledge of future Brownian increments.

If

\[
H_k
\]

is \(\mathcal{F}_{t_k}\)-measurable, then it may depend on the past and
present but not on future information.

This nonanticipation requirement is central to It\^o integration.

%%%%%%%%%%%%%%%%%%%%%%%%%%%%%%%%%%%%%%%%%%%%%%%%%%%%%%%%%%%%
\subsection{It\^o Integral for Simple Processes}
\label{subsec:ito-integral-simple}
%%%%%%%%%%%%%%%%%%%%%%%%%%%%%%%%%%%%%%%%%%%%%%%%%%%%%%%%%%%%

\begin{definitionbox}{It\^o Integral for a Simple Process}
For

\[
H_t
=
\sum_{k=0}^{n-1}
H_k
\mathbf{1}_{(t_k,t_{k+1}]}(t),
\]

define

\[
\boxed{
\int_0^T
H_t\,dB_t
=
\sum_{k=0}^{n-1}
H_k
\left(
B_{t_{k+1}}-B_{t_k}
\right).
}
\]
\end{definitionbox}

The coefficient \(H_k\) is evaluated at the left side of each interval.

%%%%%%%%%%%%%%%%%%%%%%%%%%%%%%%%%%%%%%%%%%%%%%%%%%%%%%%%%%%%
\subsection{Left-Endpoint Convention}
\label{subsec:left-endpoint-ito}
%%%%%%%%%%%%%%%%%%%%%%%%%%%%%%%%%%%%%%%%%%%%%%%%%%%%%%%%%%%%

The It\^o integral uses nonanticipating left-endpoint information.

Symbolically,

\[
\int H_t\,dB_t
\]

is built from sums resembling

\[
\boxed{
\sum_k
H_{t_k}
\left(
B_{t_{k+1}}-B_{t_k}
\right).
}
\]

This differs from ordinary Riemann integration because the integrator
\(B_t\) has unbounded variation and irregular paths.

%%%%%%%%%%%%%%%%%%%%%%%%%%%%%%%%%%%%%%%%%%%%%%%%%%%%%%%%%%%%
\subsection{Mean of an It\^o Integral}
\label{subsec:mean-ito-integral}
%%%%%%%%%%%%%%%%%%%%%%%%%%%%%%%%%%%%%%%%%%%%%%%%%%%%%%%%%%%%

For a suitable adapted integrand,

\[
\boxed{
\mathbb{E}
\left[
\int_0^T
H_t\,dB_t
\right]
=
0.
}
\]

For a simple integrand, this follows from

\[
\mathbb{E}
[
H_k
(B_{t_{k+1}}-B_{t_k})
]
=
0,
\]

because \(H_k\) depends only on the past while the Brownian increment is
independent of that past and has mean zero.

%%%%%%%%%%%%%%%%%%%%%%%%%%%%%%%%%%%%%%%%%%%%%%%%%%%%%%%%%%%%
\subsection{It\^o Isometry}
\label{subsec:ito-isometry}
%%%%%%%%%%%%%%%%%%%%%%%%%%%%%%%%%%%%%%%%%%%%%%%%%%%%%%%%%%%%

\begin{theorem}[It\^o Isometry]
For a square-integrable adapted process \(H\),

\[
\boxed{
\mathbb{E}
\left[
\left(
\int_0^T
H_t\,dB_t
\right)^2
\right]
=
\mathbb{E}
\left[
\int_0^T
H_t^2\,dt
\right].
}
\]
\end{theorem}

This identity is fundamental.

It plays a role analogous to preservation of norm under an isometric
mapping.

%%%%%%%%%%%%%%%%%%%%%%%%%%%%%%%%%%%%%%%%%%%%%%%%%%%%%%%%%%%%
\subsection{Construction by \(L^2\) Completion}
\label{subsec:ito-construction-l2}
%%%%%%%%%%%%%%%%%%%%%%%%%%%%%%%%%%%%%%%%%%%%%%%%%%%%%%%%%%%%

The It\^o integral is first defined for simple adapted processes.

More general square-integrable adapted processes are approximated by
simple processes

\[
H^{(n)}.
\]

If

\[
\mathbb{E}
\left[
\int_0^T
|H_t^{(n)}-H_t|^2\,dt
\right]
\to0,
\]

then the It\^o isometry ensures that

\[
\int_0^T
H_t^{(n)}\,dB_t
\]

is Cauchy in \(L^2\).

Its limit defines

\[
\boxed{
\int_0^T
H_t\,dB_t.
}
\]

%%%%%%%%%%%%%%%%%%%%%%%%%%%%%%%%%%%%%%%%%%%%%%%%%%%%%%%%%%%%
\subsection{Worked Example: Deterministic Integrand}
\label{subsec:worked-deterministic-ito}
%%%%%%%%%%%%%%%%%%%%%%%%%%%%%%%%%%%%%%%%%%%%%%%%%%%%%%%%%%%%

\begin{workedexamplebox}{Integral of \(t\) Against Brownian Motion}

Consider

\[
I
=
\int_0^T
t\,dB_t.
\]

The expectation is

\[
\mathbb{E}[I]=0.
\]

By the It\^o isometry,

\[
\mathbb{E}[I^2]
=
\int_0^T
t^2\,dt.
\]

Therefore,

\[
\begin{answerbox}
\[
\boxed{
\operatorname{Var}(I)
=
\frac{T^3}{3}.
}
\]
\end{answerbox}

Because the integrand is deterministic, the integral is Gaussian.

\end{workedexamplebox}

%%%%%%%%%%%%%%%%%%%%%%%%%%%%%%%%%%%%%%%%%%%%%%%%%%%%%%%%%%%%
\subsection{Gaussian It\^o Integrals with Deterministic Integrands}
\label{subsec:gaussian-ito-deterministic}
%%%%%%%%%%%%%%%%%%%%%%%%%%%%%%%%%%%%%%%%%%%%%%%%%%%%%%%%%%%%

If

\[
h:[0,T]\to\mathbb{R}
\]

is deterministic and square-integrable, then

\[
\boxed{
\int_0^T
h(t)\,dB_t
}
\]

is normally distributed with

\[
\mathbb{E}
\left[
\int_0^T
h(t)\,dB_t
\right]
=
0
\]

and

\[
\operatorname{Var}
\left(
\int_0^T
h(t)\,dB_t
\right)
=
\int_0^T
h(t)^2\,dt.
\]

Thus,

\[
\boxed{
\int_0^T
h(t)\,dB_t
\sim
N
\left(
0,
\int_0^T h(t)^2\,dt
\right).
}
\]

%%%%%%%%%%%%%%%%%%%%%%%%%%%%%%%%%%%%%%%%%%%%%%%%%%%%%%%%%%%%
\subsection{It\^o Integral as a Martingale}
\label{subsec:ito-integral-martingale}
%%%%%%%%%%%%%%%%%%%%%%%%%%%%%%%%%%%%%%%%%%%%%%%%%%%%%%%%%%%%

Define

\[
M_t
=
\int_0^t
H_s\,dB_s.
\]

Under standard square-integrability conditions,

\[
\boxed{
(M_t)_{t\geq0}
}
\]

is a martingale.

In particular,

\[
\boxed{
\mathbb{E}[M_t\mid\mathcal{F}_s]
=
M_s
}
\]

for

\[
s\leq t.
\]

%%%%%%%%%%%%%%%%%%%%%%%%%%%%%%%%%%%%%%%%%%%%%%%%%%%%%%%%%%%%
\subsection{Quadratic Variation of an It\^o Integral}
\label{subsec:qv-ito-integral}
%%%%%%%%%%%%%%%%%%%%%%%%%%%%%%%%%%%%%%%%%%%%%%%%%%%%%%%%%%%%

If

\[
M_t
=
\int_0^t
H_s\,dB_s,
\]

then its quadratic variation is

\[
\boxed{
[M]_t
=
\int_0^t
H_s^2\,ds.
}
\]

Thus the stochastic integral accumulates quadratic variation according
to the squared integrand.

%%%%%%%%%%%%%%%%%%%%%%%%%%%%%%%%%%%%%%%%%%%%%%%%%%%%%%%%%%%%
\subsection{Quadratic Covariation}
\label{subsec:quadratic-covariation}
%%%%%%%%%%%%%%%%%%%%%%%%%%%%%%%%%%%%%%%%%%%%%%%%%%%%%%%%%%%%

For suitable processes \(X\) and \(Y\), their quadratic covariation is
denoted

\[
\boxed{
[X,Y]_t.
}
\]

It is formed from limits of sums

\[
\sum_k
\Delta X_k
\Delta Y_k.
\]

Quadratic variation is the special case

\[
\boxed{
[X]_t
=
[X,X]_t.
}
\]

%%%%%%%%%%%%%%%%%%%%%%%%%%%%%%%%%%%%%%%%%%%%%%%%%%%%%%%%%%%%
\subsection{Independent Brownian Motions}
\label{subsec:independent-brownian-covariation}
%%%%%%%%%%%%%%%%%%%%%%%%%%%%%%%%%%%%%%%%%%%%%%%%%%%%%%%%%%%%

If

\[
B_t^{(1)}
\]

and

\[
B_t^{(2)}
\]

are independent Brownian motions, then

\[
\boxed{
[B^{(1)},B^{(2)}]_t
=
0.
}
\]

Symbolically, one often writes

\[
\boxed{
dB_t^{(1)}
dB_t^{(2)}
=
0.
}
\]

%%%%%%%%%%%%%%%%%%%%%%%%%%%%%%%%%%%%%%%%%%%%%%%%%%%%%%%%%%%%
\subsection{Correlated Brownian Motions}
\label{subsec:correlated-brownian-motions}
%%%%%%%%%%%%%%%%%%%%%%%%%%%%%%%%%%%%%%%%%%%%%%%%%%%%%%%%%%%%

Suppose Brownian motions satisfy instantaneous correlation \(\rho\).

Then their quadratic covariation satisfies

\[
\boxed{
d[B^{(1)},B^{(2)}]_t
=
\rho\,dt.
}
\]

Symbolically,

\[
\boxed{
dB_t^{(1)}
dB_t^{(2)}
=
\rho\,dt.
}
\]

This relation is widely used in multivariate diffusion models.

%%%%%%%%%%%%%%%%%%%%%%%%%%%%%%%%%%%%%%%%%%%%%%%%%%%%%%%%%%%%
\subsection{It\^o's Formula}
\label{subsec:ito-formula}
%%%%%%%%%%%%%%%%%%%%%%%%%%%%%%%%%%%%%%%%%%%%%%%%%%%%%%%%%%%%

It\^o's formula is the stochastic analogue of the chain rule.

\begin{theorem}[It\^o's Formula]
Let

\[
X_t
\]

satisfy

\[
dX_t
=
\mu_t\,dt
+
\sigma_t\,dB_t,
\]

and let

\[
f(t,x)
\]

have suitable derivatives.

Then

\[
\boxed{
\begin{aligned}
df(t,X_t)
={}&
\frac{\partial f}{\partial t}(t,X_t)\,dt
+
\frac{\partial f}{\partial x}(t,X_t)\,dX_t
\\
&+
\frac12
\frac{\partial^2 f}{\partial x^2}(t,X_t)
\sigma_t^2\,dt.
\end{aligned}
}
\]
\end{theorem}

Substituting \(dX_t\),

\[
\boxed{
\begin{aligned}
df(t,X_t)
={}&
\left(
f_t
+
\mu_t f_x
+
\frac12\sigma_t^2f_{xx}
\right)dt
\\
&+
\sigma_t f_x\,dB_t.
\end{aligned}
}
\]

%%%%%%%%%%%%%%%%%%%%%%%%%%%%%%%%%%%%%%%%%%%%%%%%%%%%%%%%%%%%
\subsection{Why It\^o's Formula Has a Second Derivative}
\label{subsec:why-ito-second-derivative}
%%%%%%%%%%%%%%%%%%%%%%%%%%%%%%%%%%%%%%%%%%%%%%%%%%%%%%%%%%%%

Ordinary Taylor expansion gives

\[
df
\approx
f_x\,dX
+
\frac12
f_{xx}(dX)^2.
\]

If

\[
dX
=
\mu\,dt+\sigma\,dB,
\]

then

\[
(dX)^2
\]

contains the term

\[
\sigma^2(dB)^2.
\]

Using

\[
(dB)^2=dt,
\]

we obtain

\[
(dX)^2
=
\sigma^2dt
\]

to first stochastic order.

Therefore the second derivative term survives.

%%%%%%%%%%%%%%%%%%%%%%%%%%%%%%%%%%%%%%%%%%%%%%%%%%%%%%%%%%%%
\subsection{It\^o Multiplication Table}
\label{subsec:ito-multiplication-table}
%%%%%%%%%%%%%%%%%%%%%%%%%%%%%%%%%%%%%%%%%%%%%%%%%%%%%%%%%%%%

The symbolic rules are summarized by

\[
\begin{array}{c|cc}
 & dt & dB_t\\
\hline
dt & 0 & 0\\
dB_t & 0 & dt
\end{array}
\]

meaning

\[
\boxed{
(dt)^2=0,
}
\]

\[
\boxed{
dt\,dB_t=0,
}
\]

and

\[
\boxed{
(dB_t)^2=dt.
}
\]

%%%%%%%%%%%%%%%%%%%%%%%%%%%%%%%%%%%%%%%%%%%%%%%%%%%%%%%%%%%%
\subsection{Worked Example: \(B_t^2\)}
\label{subsec:worked-ito-brownian-square}
%%%%%%%%%%%%%%%%%%%%%%%%%%%%%%%%%%%%%%%%%%%%%%%%%%%%%%%%%%%%

\begin{workedexamplebox}{Differential of \(B_t^2\)}

Let

\[
f(x)=x^2.
\]

Then

\[
f'(x)=2x,
\qquad
f''(x)=2.
\]

Since

\[
dB_t=dB_t,
\]

It\^o's formula gives

\[
d(B_t^2)
=
2B_t\,dB_t
+
\frac12(2)\,dt.
\]

Therefore,

\begin{answerbox}
\[
\boxed{
d(B_t^2)
=
2B_t\,dB_t
+
dt.
}
\]
\end{answerbox}

Integrating,

\[
\boxed{
B_t^2
=
2\int_0^t
B_s\,dB_s
+
t.
}
\]

\end{workedexamplebox}

%%%%%%%%%%%%%%%%%%%%%%%%%%%%%%%%%%%%%%%%%%%%%%%%%%%%%%%%%%%%
\subsection{Computing \(\int B_t\,dB_t\)}
\label{subsec:integral-b-db}
%%%%%%%%%%%%%%%%%%%%%%%%%%%%%%%%%%%%%%%%%%%%%%%%%%%%%%%%%%%%

From

\[
B_t^2
=
2\int_0^t
B_s\,dB_s+t,
\]

we obtain

\[
\boxed{
\int_0^t
B_s\,dB_s
=
\frac12
(B_t^2-t).
}
\]

This differs from the ordinary-calculus expression

\[
\frac12B_t^2
\]

by the correction term

\[
-\frac{t}{2}.
\]

%%%%%%%%%%%%%%%%%%%%%%%%%%%%%%%%%%%%%%%%%%%%%%%%%%%%%%%%%%%%
\subsection{Worked Example: Expectation of Brownian Square}
\label{subsec:worked-expectation-brownian-square}
%%%%%%%%%%%%%%%%%%%%%%%%%%%%%%%%%%%%%%%%%%%%%%%%%%%%%%%%%%%%

\begin{workedexamplebox}{Recovering Brownian Variance}

Take expectations in

\[
B_t^2
=
2\int_0^tB_s\,dB_s+t.
\]

The It\^o integral has mean zero.

Therefore,

\[
\mathbb{E}[B_t^2]
=
t.
\]

\begin{answerbox}
\[
\boxed{
\operatorname{Var}(B_t)=t.
}
\]
\end{answerbox}

\end{workedexamplebox}

%%%%%%%%%%%%%%%%%%%%%%%%%%%%%%%%%%%%%%%%%%%%%%%%%%%%%%%%%%%%
\subsection{Worked Example: Exponential of Brownian Motion}
\label{subsec:worked-ito-exp-brownian}
%%%%%%%%%%%%%%%%%%%%%%%%%%%%%%%%%%%%%%%%%%%%%%%%%%%%%%%%%%%%

\begin{workedexamplebox}{Differential of \(e^{B_t}\)}

Let

\[
f(x)=e^x.
\]

Then

\[
f'(x)=e^x
\]

and

\[
f''(x)=e^x.
\]

It\^o's formula gives

\[
d(e^{B_t})
=
e^{B_t}\,dB_t
+
\frac12e^{B_t}\,dt.
\]

\begin{answerbox}
\[
\boxed{
d(e^{B_t})
=
\frac12e^{B_t}\,dt
+
e^{B_t}\,dB_t.
}
\]
\end{answerbox}

\end{workedexamplebox}

%%%%%%%%%%%%%%%%%%%%%%%%%%%%%%%%%%%%%%%%%%%%%%%%%%%%%%%%%%%%
\subsection{Exponential Martingale}
\label{subsec:exponential-martingale}
%%%%%%%%%%%%%%%%%%%%%%%%%%%%%%%%%%%%%%%%%%%%%%%%%%%%%%%%%%%%

Consider

\[
\boxed{
M_t
=
\exp
\left(
\theta B_t
-
\frac12\theta^2t
\right).
}
\]

Apply It\^o's formula.

The drift correction

\[
-\frac12\theta^2t
\]

exactly cancels the It\^o second-order term.

One obtains

\[
\boxed{
dM_t
=
\theta M_t\,dB_t.
}
\]

Thus, under standard conditions,

\[
\boxed{
(M_t)_{t\geq0}
}
\]

is a martingale.

%%%%%%%%%%%%%%%%%%%%%%%%%%%%%%%%%%%%%%%%%%%%%%%%%%%%%%%%%%%%
\subsection{Worked Example: Why the Drift Correction Appears}
\label{subsec:worked-exponential-martingale}
%%%%%%%%%%%%%%%%%%%%%%%%%%%%%%%%%%%%%%%%%%%%%%%%%%%%%%%%%%%%

\begin{workedexamplebox}{Constructing a Martingale}

Let

\[
M_t
=
e^{B_t-t/2}.
\]

Set

\[
f(t,x)
=
e^{x-t/2}.
\]

Then

\[
f_t
=
-\frac12f,
\]

\[
f_x=f,
\]

and

\[
f_{xx}=f.
\]

It\^o's formula gives

\[
dM_t
=
\left(
-\frac12f+\frac12f
\right)dt
+
f\,dB_t.
\]

The drift terms cancel.

\begin{answerbox}
\[
\boxed{
dM_t
=
M_t\,dB_t.
}
\]
\end{answerbox}

Hence \(M_t\) is a martingale under the standard integrability
conditions.

\end{workedexamplebox}

%%%%%%%%%%%%%%%%%%%%%%%%%%%%%%%%%%%%%%%%%%%%%%%%%%%%%%%%%%%%
\subsection{Time-Dependent It\^o Formula}
\label{subsec:time-dependent-ito}
%%%%%%%%%%%%%%%%%%%%%%%%%%%%%%%%%%%%%%%%%%%%%%%%%%%%%%%%%%%%

For

\[
f=f(t,B_t),
\]

It\^o's formula becomes

\[
\boxed{
df
=
f_t\,dt
+
f_x\,dB_t
+
\frac12f_{xx}\,dt.
}
\]

Equivalently,

\[
\boxed{
df
=
\left(
f_t+\frac12f_{xx}
\right)dt
+
f_x\,dB_t.
}
\]

The differential operator

\[
\boxed{
\frac{\partial}{\partial t}
+
\frac12
\frac{\partial^2}{\partial x^2}
}
\]

appears naturally.

%%%%%%%%%%%%%%%%%%%%%%%%%%%%%%%%%%%%%%%%%%%%%%%%%%%%%%%%%%%%
\subsection{Connection to the Heat Equation}
\label{subsec:ito-heat-equation}
%%%%%%%%%%%%%%%%%%%%%%%%%%%%%%%%%%%%%%%%%%%%%%%%%%%%%%%%%%%%

Suppose

\[
f_t+\frac12f_{xx}=0.
\]

Then

\[
df(t,B_t)
=
f_x(t,B_t)\,dB_t.
\]

The drift vanishes.

Therefore

\[
f(t,B_t)
\]

is a local martingale under suitable conditions.

This connects Brownian motion directly with the heat equation.

%%%%%%%%%%%%%%%%%%%%%%%%%%%%%%%%%%%%%%%%%%%%%%%%%%%%%%%%%%%%
\subsection{Generator of Brownian Motion}
\label{subsec:brownian-generator}
%%%%%%%%%%%%%%%%%%%%%%%%%%%%%%%%%%%%%%%%%%%%%%%%%%%%%%%%%%%%

The infinitesimal generator of standard Brownian motion is

\[
\boxed{
\mathcal{L}
=
\frac12
\frac{d^2}{dx^2}.
}
\]

For a diffusion with

\[
dX_t
=
\mu(X_t,t)\,dt
+
\sigma(X_t,t)\,dB_t,
\]

the corresponding formal generator is

\[
\boxed{
\mathcal{L}f
=
\mu f_x
+
\frac12\sigma^2f_{xx}.
}
\]

This operator connects stochastic processes with partial differential
equations.

%%%%%%%%%%%%%%%%%%%%%%%%%%%%%%%%%%%%%%%%%%%%%%%%%%%%%%%%%%%%
\subsection{Dynkin's Formula Preview}
\label{subsec:dynkin-formula-preview}
%%%%%%%%%%%%%%%%%%%%%%%%%%%%%%%%%%%%%%%%%%%%%%%%%%%%%%%%%%%%

For a suitable Markov diffusion with generator \(\mathcal{L}\),

\[
\boxed{
\mathbb{E}_x[f(X_t)]
=
f(x)
+
\mathbb{E}_x
\left[
\int_0^t
\mathcal{L}f(X_s)\,ds
\right].
}
\]

This is Dynkin's formula.

It can be viewed as an expectation-level consequence of It\^o's
formula.

%%%%%%%%%%%%%%%%%%%%%%%%%%%%%%%%%%%%%%%%%%%%%%%%%%%%%%%%%%%%
\subsection{Stochastic Product Rule}
\label{subsec:stochastic-product-rule}
%%%%%%%%%%%%%%%%%%%%%%%%%%%%%%%%%%%%%%%%%%%%%%%%%%%%%%%%%%%%

For suitable semimartingales \(X_t\) and \(Y_t\),

\[
\boxed{
d(X_tY_t)
=
X_t\,dY_t
+
Y_t\,dX_t
+
d[X,Y]_t.
}
\]

The final quadratic-covariation term is the stochastic correction.

In ordinary smooth calculus, that term vanishes.

%%%%%%%%%%%%%%%%%%%%%%%%%%%%%%%%%%%%%%%%%%%%%%%%%%%%%%%%%%%%
\subsection{Brownian Product Rule}
\label{subsec:brownian-product-rule}
%%%%%%%%%%%%%%%%%%%%%%%%%%%%%%%%%%%%%%%%%%%%%%%%%%%%%%%%%%%%

For two Brownian motions,

\[
d(B_t^{(1)}B_t^{(2)})
=
B_t^{(1)}\,dB_t^{(2)}
+
B_t^{(2)}\,dB_t^{(1)}
+
d[B^{(1)},B^{(2)}]_t.
\]

If they are independent,

\[
d[B^{(1)},B^{(2)}]_t=0.
\]

If their instantaneous correlation is \(\rho\),

\[
d[B^{(1)},B^{(2)}]_t
=
\rho\,dt.
\]

%%%%%%%%%%%%%%%%%%%%%%%%%%%%%%%%%%%%%%%%%%%%%%%%%%%%%%%%%%%%
\subsection{Stochastic Integration by Parts}
\label{subsec:stochastic-integration-parts}
%%%%%%%%%%%%%%%%%%%%%%%%%%%%%%%%%%%%%%%%%%%%%%%%%%%%%%%%%%%%

Integrating the stochastic product rule gives

\[
\boxed{
X_tY_t
=
X_0Y_0
+
\int_0^t
X_s\,dY_s
+
\int_0^t
Y_s\,dX_s
+
[X,Y]_t.
}
\]

Thus stochastic integration by parts contains an additional
quadratic-covariation term.

%%%%%%%%%%%%%%%%%%%%%%%%%%%%%%%%%%%%%%%%%%%%%%%%%%%%%%%%%%%%
\subsection{Worked Example: Integration by Parts with \(tB_t\)}
\label{subsec:worked-integration-parts-tb}
%%%%%%%%%%%%%%%%%%%%%%%%%%%%%%%%%%%%%%%%%%%%%%%%%%%%%%%%%%%%

\begin{workedexamplebox}{Computing \(\int_0^t s\,dB_s\)}

Take

\[
X_t=t,
\qquad
Y_t=B_t.
\]

Since \(t\) has zero quadratic variation with Brownian motion,

\[
[t,B]_t=0.
\]

The product rule gives

\[
d(tB_t)
=
t\,dB_t
+
B_t\,dt.
\]

Integrating,

\[
tB_t
=
\int_0^t
s\,dB_s
+
\int_0^t
B_s\,ds.
\]

Therefore,

\begin{answerbox}
\[
\boxed{
\int_0^t
s\,dB_s
=
tB_t
-
\int_0^t
B_s\,ds.
}
\]
\end{answerbox}

\end{workedexamplebox}

%%%%%%%%%%%%%%%%%%%%%%%%%%%%%%%%%%%%%%%%%%%%%%%%%%%%%%%%%%%%
\subsection{Multidimensional It\^o Formula}
\label{subsec:multidimensional-ito}
%%%%%%%%%%%%%%%%%%%%%%%%%%%%%%%%%%%%%%%%%%%%%%%%%%%%%%%%%%%%

Suppose

\[
\mathbf{X}_t
=
(X_t^{(1)},\ldots,X_t^{(n)})
\]

is a vector It\^o process.

For a smooth function

\[
f(t,\mathbf{x}),
\]

the multidimensional It\^o formula contains:

\[
\boxed{
\text{time derivative}
}
\]

\[
\boxed{
\text{first spatial derivatives}
}
\]

and

\[
\boxed{
\text{second derivatives weighted by quadratic covariations}.
}
\]

In differential notation,

\[
\boxed{
df
=
f_t\,dt
+
\sum_i
f_{x_i}\,dX_t^{(i)}
+
\frac12
\sum_{i,j}
f_{x_ix_j}
\,d[X^{(i)},X^{(j)}]_t.
}
\]

%%%%%%%%%%%%%%%%%%%%%%%%%%%%%%%%%%%%%%%%%%%%%%%%%%%%%%%%%%%%
\subsection{Multidimensional Brownian Motion}
\label{subsec:multidimensional-brownian}
%%%%%%%%%%%%%%%%%%%%%%%%%%%%%%%%%%%%%%%%%%%%%%%%%%%%%%%%%%%%

A standard \(d\)-dimensional Brownian motion is

\[
\boxed{
\mathbf{B}_t
=
(B_t^{(1)},\ldots,B_t^{(d)}),
}
\]

where the coordinate Brownian motions are independent.

Then

\[
\boxed{
dB_t^{(i)}
dB_t^{(j)}
=
\delta_{ij}\,dt.
}
\]

Here

\[
\delta_{ij}
\]

is the Kronecker delta:

\[
\delta_{ij}
=
\begin{cases}
1,&i=j,\\
0,&i\neq j.
\end{cases}
\]

%%%%%%%%%%%%%%%%%%%%%%%%%%%%%%%%%%%%%%%%%%%%%%%%%%%%%%%%%%%%
\subsection{Worked Example: Squared Radius of Brownian Motion}
\label{subsec:worked-brownian-radius-square}
%%%%%%%%%%%%%%%%%%%%%%%%%%%%%%%%%%%%%%%%%%%%%%%%%%%%%%%%%%%%

\begin{workedexamplebox}{Two-Dimensional Brownian Radius}

Let

\[
R_t^2
=
(B_t^{(1)})^2
+
(B_t^{(2)})^2.
\]

For each coordinate,

\[
d(B_t^{(i)})^2
=
2B_t^{(i)}\,dB_t^{(i)}
+
dt.
\]

Adding,

\[
dR_t^2
=
2B_t^{(1)}\,dB_t^{(1)}
+
2B_t^{(2)}\,dB_t^{(2)}
+
2dt.
\]

\begin{answerbox}
\[
\boxed{
dR_t^2
=
2\mathbf{B}_t\cdot d\mathbf{B}_t
+
2dt.
}
\]
\end{answerbox}

\end{workedexamplebox}

%%%%%%%%%%%%%%%%%%%%%%%%%%%%%%%%%%%%%%%%%%%%%%%%%%%%%%%%%%%%
\subsection{It\^o Processes}
\label{subsec:ito-processes}
%%%%%%%%%%%%%%%%%%%%%%%%%%%%%%%%%%%%%%%%%%%%%%%%%%%%%%%%%%%%

A process of the form

\[
\boxed{
X_t
=
X_0
+
\int_0^t
\mu_s\,ds
+
\int_0^t
\sigma_s\,dB_s
}
\]

is called an It\^o process under suitable measurability and integrability
conditions.

In differential notation,

\[
\boxed{
dX_t
=
\mu_t\,dt
+
\sigma_t\,dB_t.
}
\]

Here:

\[
\mu_t
\]

is the drift,

and

\[
\sigma_t
\]

is the diffusion coefficient or volatility.

%%%%%%%%%%%%%%%%%%%%%%%%%%%%%%%%%%%%%%%%%%%%%%%%%%%%%%%%%%%%
\subsection{Drift and Diffusion}
\label{subsec:drift-diffusion}
%%%%%%%%%%%%%%%%%%%%%%%%%%%%%%%%%%%%%%%%%%%%%%%%%%%%%%%%%%%%

In

\[
dX_t
=
\mu_t\,dt
+
\sigma_t\,dB_t,
\]

the term

\[
\mu_t\,dt
\]

describes predictable local motion.

The term

\[
\sigma_t\,dB_t
\]

describes random local fluctuation.

Thus:

\[
\boxed{
\text{It\^o process}
=
\text{drift}
+
\text{diffusion noise}.
}
\]

%%%%%%%%%%%%%%%%%%%%%%%%%%%%%%%%%%%%%%%%%%%%%%%%%%%%%%%%%%%%
\subsection{Quadratic Variation of an It\^o Process}
\label{subsec:qv-ito-process}
%%%%%%%%%%%%%%%%%%%%%%%%%%%%%%%%%%%%%%%%%%%%%%%%%%%%%%%%%%%%

For

\[
dX_t
=
\mu_t\,dt
+
\sigma_t\,dB_t,
\]

the finite-variation drift contributes no quadratic variation.

Therefore,

\[
\boxed{
d[X]_t
=
\sigma_t^2\,dt.
}
\]

Equivalently,

\[
\boxed{
[X]_t
=
\int_0^t
\sigma_s^2\,ds.
}
\]

%%%%%%%%%%%%%%%%%%%%%%%%%%%%%%%%%%%%%%%%%%%%%%%%%%%%%%%%%%%%
\subsection{Local Martingales}
\label{subsec:local-martingales}
%%%%%%%%%%%%%%%%%%%%%%%%%%%%%%%%%%%%%%%%%%%%%%%%%%%%%%%%%%%%

A local martingale is a process that becomes a martingale after
stopping at an increasing sequence of stopping times.

Roughly, there exist stopping times

\[
\tau_n\uparrow\infty
\]

such that

\[
\boxed{
M_{t\wedge\tau_n}
}
\]

is a martingale for every \(n\).

Local martingales are more general than true martingales.

%%%%%%%%%%%%%%%%%%%%%%%%%%%%%%%%%%%%%%%%%%%%%%%%%%%%%%%%%%%%
\subsection{Why Local Martingales Matter}
\label{subsec:why-local-martingales}
%%%%%%%%%%%%%%%%%%%%%%%%%%%%%%%%%%%%%%%%%%%%%%%%%%%%%%%%%%%%

Many stochastic-calculus expressions naturally produce local
martingales.

Additional integrability conditions may be needed to conclude that they
are true martingales.

Thus one should not automatically replace

\[
\boxed{
\text{local martingale}
}
\]

with

\[
\boxed{
\text{martingale}.
}
\]

%%%%%%%%%%%%%%%%%%%%%%%%%%%%%%%%%%%%%%%%%%%%%%%%%%%%%%%%%%%%
\subsection{Stochastic Exponential}
\label{subsec:stochastic-exponential}
%%%%%%%%%%%%%%%%%%%%%%%%%%%%%%%%%%%%%%%%%%%%%%%%%%%%%%%%%%%%

For a continuous local martingale \(M_t\), the stochastic exponential is

\[
\boxed{
\mathcal{E}(M)_t
=
\exp
\left(
M_t
-
\frac12[M]_t
\right).
}
\]

For

\[
M_t
=
\theta B_t,
\]

this becomes

\[
\boxed{
\mathcal{E}(\theta B)_t
=
\exp
\left(
\theta B_t
-
\frac12\theta^2t
\right).
}
\]

This recovers the Brownian exponential martingale.

%%%%%%%%%%%%%%%%%%%%%%%%%%%%%%%%%%%%%%%%%%%%%%%%%%%%%%%%%%%%
\subsection{Novikov Condition Preview}
\label{subsec:novikov-preview}
%%%%%%%%%%%%%%%%%%%%%%%%%%%%%%%%%%%%%%%%%%%%%%%%%%%%%%%%%%%%

A local stochastic exponential need not always be a true martingale.

A useful sufficient condition is Novikov's condition.

For a process \(\theta_t\), a standard form is

\[
\boxed{
\mathbb{E}
\left[
\exp
\left(
\frac12
\int_0^T
\theta_t^2\,dt
\right)
\right]
<
\infty.
}
\]

Under suitable conditions, the associated stochastic exponential is a
true martingale.

%%%%%%%%%%%%%%%%%%%%%%%%%%%%%%%%%%%%%%%%%%%%%%%%%%%%%%%%%%%%
\subsection{Stratonovich Integral}
\label{subsec:stratonovich-integral}
%%%%%%%%%%%%%%%%%%%%%%%%%%%%%%%%%%%%%%%%%%%%%%%%%%%%%%%%%%%%

Another stochastic integral is the Stratonovich integral, written

\[
\boxed{
\int_0^t
H_s
\circ dB_s.
}
\]

It is based conceptually on midpoint-style approximations rather than
the left-endpoint convention of It\^o integration.

Its notation contains a small circle:

\[
\circ dB_s.
\]

%%%%%%%%%%%%%%%%%%%%%%%%%%%%%%%%%%%%%%%%%%%%%%%%%%%%%%%%%%%%
\subsection{It\^o Versus Stratonovich}
\label{subsec:ito-versus-stratonovich}
%%%%%%%%%%%%%%%%%%%%%%%%%%%%%%%%%%%%%%%%%%%%%%%%%%%%%%%%%%%%

The It\^o integral has especially strong martingale properties.

The Stratonovich integral behaves more like ordinary calculus under
smooth coordinate changes.

For a sufficiently regular process \(H_t\),

\[
\boxed{
\int_0^t
H_s
\circ dB_s
=
\int_0^t
H_s\,dB_s
+
\frac12[H,B]_t.
}
\]

Thus the two integrals differ by a quadratic-covariation correction.

%%%%%%%%%%%%%%%%%%%%%%%%%%%%%%%%%%%%%%%%%%%%%%%%%%%%%%%%%%%%
\subsection{Worked Example: It\^o Versus Stratonovich}
\label{subsec:worked-ito-stratonovich}
%%%%%%%%%%%%%%%%%%%%%%%%%%%%%%%%%%%%%%%%%%%%%%%%%%%%%%%%%%%%

\begin{workedexamplebox}{Integrating Brownian Motion Against Itself}

For It\^o integration,

\[
\int_0^t
B_s\,dB_s
=
\frac12
(B_t^2-t).
\]

For Stratonovich integration, ordinary chain-rule behavior gives

\[
\int_0^t
B_s\circ dB_s
=
\frac12B_t^2.
\]

Therefore,

\begin{answerbox}
\[
\boxed{
\int_0^t
B_s\circ dB_s
=
\int_0^t
B_s\,dB_s
+
\frac12t.
}
\]
\end{answerbox}

\end{workedexamplebox}

%%%%%%%%%%%%%%%%%%%%%%%%%%%%%%%%%%%%%%%%%%%%%%%%%%%%%%%%%%%%
\subsection{Stratonovich Chain Rule}
\label{subsec:stratonovich-chain-rule}
%%%%%%%%%%%%%%%%%%%%%%%%%%%%%%%%%%%%%%%%%%%%%%%%%%%%%%%%%%%%

For sufficiently regular processes, Stratonovich calculus formally
obeys the ordinary chain rule:

\[
\boxed{
df(X_t)
=
f'(X_t)
\circ dX_t.
}
\]

The It\^o formulation instead introduces the second-order correction.

This is one reason Stratonovich notation often appears in physics and
geometric stochastic analysis.

%%%%%%%%%%%%%%%%%%%%%%%%%%%%%%%%%%%%%%%%%%%%%%%%%%%%%%%%%%%%
\subsection{Elementary It\^o--Stratonovich Conversion}
\label{subsec:ito-stratonovich-conversion}
%%%%%%%%%%%%%%%%%%%%%%%%%%%%%%%%%%%%%%%%%%%%%%%%%%%%%%%%%%%%

Consider the Stratonovich equation

\[
dX_t
=
a(X_t)\,dt
+
b(X_t)\circ dB_t.
\]

The corresponding one-dimensional It\^o form is

\[
\boxed{
dX_t
=
\left[
a(X_t)
+
\frac12
b(X_t)b'(X_t)
\right]dt
+
b(X_t)\,dB_t.
}
\]

Thus changing stochastic conventions changes the drift term.

%%%%%%%%%%%%%%%%%%%%%%%%%%%%%%%%%%%%%%%%%%%%%%%%%%%%%%%%%%%%
\subsection{Girsanov's Theorem Preview}
\label{subsec:girsanov-preview}
%%%%%%%%%%%%%%%%%%%%%%%%%%%%%%%%%%%%%%%%%%%%%%%%%%%%%%%%%%%%

Girsanov's theorem describes how Brownian drift changes under an
absolutely continuous change of probability measure.

Very roughly, suppose a process has the form

\[
dX_t
=
\mu_t\,dt
+
\sigma_t\,dB_t.
\]

Under a suitably chosen new probability measure, one may absorb part of
the drift into a transformed Brownian motion.

This principle is fundamental in mathematical finance and stochastic
control.

%%%%%%%%%%%%%%%%%%%%%%%%%%%%%%%%%%%%%%%%%%%%%%%%%%%%%%%%%%%%
\subsection{Change of Measure}
\label{subsec:change-of-measure-stochastic}
%%%%%%%%%%%%%%%%%%%%%%%%%%%%%%%%%%%%%%%%%%%%%%%%%%%%%%%%%%%%

Suppose \(Q\) is another probability measure satisfying

\[
Q\ll P.
\]

The Radon--Nikodym derivative

\[
\boxed{
\frac{dQ}{dP}
}
\]

describes the change from \(P\) to \(Q\).

In stochastic calculus, stochastic exponentials often serve as density
processes for such measure changes.

%%%%%%%%%%%%%%%%%%%%%%%%%%%%%%%%%%%%%%%%%%%%%%%%%%%%%%%%%%%%
\subsection{Girsanov Density Process Preview}
\label{subsec:girsanov-density-preview}
%%%%%%%%%%%%%%%%%%%%%%%%%%%%%%%%%%%%%%%%%%%%%%%%%%%%%%%%%%%%

A typical stochastic exponential is

\[
\boxed{
Z_t
=
\exp
\left(
-\int_0^t
\theta_s\,dB_s
-
\frac12
\int_0^t
\theta_s^2\,ds
\right).
}
\]

Under suitable conditions,

\[
Z_t
\]

is a martingale and can define a new measure.

Under that measure, the process

\[
\boxed{
W_t
=
B_t
+
\int_0^t
\theta_s\,ds
}
\]

can become Brownian motion.

%%%%%%%%%%%%%%%%%%%%%%%%%%%%%%%%%%%%%%%%%%%%%%%%%%%%%%%%%%%%
\subsection{Feynman--Kac Formula Preview}
\label{subsec:feynman-kac-preview}
%%%%%%%%%%%%%%%%%%%%%%%%%%%%%%%%%%%%%%%%%%%%%%%%%%%%%%%%%%%%

Stochastic calculus creates deep connections between diffusion
processes and partial differential equations.

A Feynman--Kac formula represents solutions of certain PDEs as
expectations of stochastic processes.

A typical expression has the conceptual form

\[
\boxed{
u(t,x)
=
\mathbb{E}_{t,x}
\left[
\text{discounted future payoff}
\right].
}
\]

This connection is central in mathematical finance, physics, and
stochastic control.

%%%%%%%%%%%%%%%%%%%%%%%%%%%%%%%%%%%%%%%%%%%%%%%%%%%%%%%%%%%%
\subsection{Stochastic Calculus and PDEs}
\label{subsec:stochastic-calculus-pdes}
%%%%%%%%%%%%%%%%%%%%%%%%%%%%%%%%%%%%%%%%%%%%%%%%%%%%%%%%%%%%

The generator

\[
\mathcal{L}
=
\mu\frac{\partial}{\partial x}
+
\frac12\sigma^2
\frac{\partial^2}{\partial x^2}
\]

naturally appears in PDEs such as

\[
\boxed{
u_t+\mathcal{L}u=0.
}
\]

The second-order derivative arises directly from Brownian quadratic
variation.

Thus stochastic calculus and second-order PDEs are structurally linked.

%%%%%%%%%%%%%%%%%%%%%%%%%%%%%%%%%%%%%%%%%%%%%%%%%%%%%%%%%%%%
\subsection{Stochastic Calculus in Finance}
\label{subsec:stochastic-calculus-finance}
%%%%%%%%%%%%%%%%%%%%%%%%%%%%%%%%%%%%%%%%%%%%%%%%%%%%%%%%%%%%

A stylized asset model may be written

\[
\boxed{
dS_t
=
\mu S_t\,dt
+
\sigma S_t\,dB_t.
}
\]

The first term represents expected growth.

The second represents random proportional fluctuations.

Applying It\^o's formula to

\[
\ln S_t
\]

produces a correction term involving

\[
-\frac12\sigma^2.
\]

The full solution is developed in the next section on stochastic
differential equations.

%%%%%%%%%%%%%%%%%%%%%%%%%%%%%%%%%%%%%%%%%%%%%%%%%%%%%%%%%%%%
\subsection{Worked Example: Differential of \(\ln X_t\)}
\label{subsec:worked-ito-log}
%%%%%%%%%%%%%%%%%%%%%%%%%%%%%%%%%%%%%%%%%%%%%%%%%%%%%%%%%%%%

\begin{workedexamplebox}{Logarithm of an It\^o Process}

Suppose

\[
dX_t
=
\mu X_t\,dt
+
\sigma X_t\,dB_t,
\qquad
X_t>0.
\]

Let

\[
f(x)=\ln x.
\]

Then

\[
f'(x)=\frac1x,
\]

and

\[
f''(x)
=
-\frac1{x^2}.
\]

It\^o's formula gives

\[
d(\ln X_t)
=
\frac1{X_t}dX_t
-
\frac12
\frac1{X_t^2}
\sigma^2X_t^2\,dt.
\]

Therefore,

\begin{answerbox}
\[
\boxed{
d(\ln X_t)
=
\left(
\mu-\frac12\sigma^2
\right)dt
+
\sigma\,dB_t.
}
\]
\end{answerbox}

\end{workedexamplebox}

%%%%%%%%%%%%%%%%%%%%%%%%%%%%%%%%%%%%%%%%%%%%%%%%%%%%%%%%%%%%
\subsection{It\^o Formula and Expected Values}
\label{subsec:ito-expected-values}
%%%%%%%%%%%%%%%%%%%%%%%%%%%%%%%%%%%%%%%%%%%%%%%%%%%%%%%%%%%%

Suppose

\[
dX_t
=
\mu_t\,dt
+
\sigma_t\,dB_t
\]

and \(f\) is sufficiently regular.

It\^o's formula gives

\[
df(X_t)
=
\left(
\mu_tf'(X_t)
+
\frac12\sigma_t^2f''(X_t)
\right)dt
+
\sigma_tf'(X_t)\,dB_t.
\]

Taking expectations and using the zero mean of the It\^o integral gives

\[
\boxed{
\frac{d}{dt}
\mathbb{E}[f(X_t)]
=
\mathbb{E}
\left[
\mu_tf'(X_t)
+
\frac12\sigma_t^2f''(X_t)
\right]
}
\]

under appropriate conditions.

%%%%%%%%%%%%%%%%%%%%%%%%%%%%%%%%%%%%%%%%%%%%%%%%%%%%%%%%%%%%
\subsection{Occupation Time Preview}
\label{subsec:occupation-time-preview}
%%%%%%%%%%%%%%%%%%%%%%%%%%%%%%%%%%%%%%%%%%%%%%%%%%%%%%%%%%%%

Stochastic processes may spend random amounts of time in different
regions.

For a process \(X_t\), an occupation time of a set \(A\) is

\[
\boxed{
\int_0^T
\mathbf{1}_{\{X_t\in A\}}
\,dt.
}
\]

Brownian occupation times lead to the concept of local time.

%%%%%%%%%%%%%%%%%%%%%%%%%%%%%%%%%%%%%%%%%%%%%%%%%%%%%%%%%%%%
\subsection{Local Time Preview}
\label{subsec:local-time-preview}
%%%%%%%%%%%%%%%%%%%%%%%%%%%%%%%%%%%%%%%%%%%%%%%%%%%%%%%%%%%%

Brownian local time

\[
L_t^a
\]

measures, in an appropriate density sense, how much time Brownian motion
spends near a level \(a\).

A central identity involving local time is Tanaka's formula.

Local time becomes important in:

\[
\text{reflected diffusions},
\]

\[
\text{optimal stopping},
\]

and

\[
\text{nonsmooth stochastic calculus}.
\]

%%%%%%%%%%%%%%%%%%%%%%%%%%%%%%%%%%%%%%%%%%%%%%%%%%%%%%%%%%%%
\subsection{Tanaka's Formula Preview}
\label{subsec:tanaka-preview}
%%%%%%%%%%%%%%%%%%%%%%%%%%%%%%%%%%%%%%%%%%%%%%%%%%%%%%%%%%%%

Ordinary It\^o's formula requires a second derivative.

The function

\[
|x|
\]

is not twice differentiable at \(0\).

Tanaka's formula extends stochastic calculus to such nonsmooth
functions and introduces local time:

\[
\boxed{
|B_t|
=
|B_0|
+
\int_0^t
\operatorname{sgn}(B_s)\,dB_s
+
L_t^0.
}
\]

This reveals how local time compensates for the nondifferentiability at
the origin.

%%%%%%%%%%%%%%%%%%%%%%%%%%%%%%%%%%%%%%%%%%%%%%%%%%%%%%%%%%%%
\subsection{Semimartingales Preview}
\label{subsec:semimartingales-preview}
%%%%%%%%%%%%%%%%%%%%%%%%%%%%%%%%%%%%%%%%%%%%%%%%%%%%%%%%%%%%

The broad class of processes suitable for It\^o-style stochastic
integration is the class of semimartingales.

A semimartingale can be decomposed conceptually as

\[
\boxed{
\text{local martingale}
+
\text{finite-variation process}.
}
\]

It\^o processes are important examples of semimartingales.

%%%%%%%%%%%%%%%%%%%%%%%%%%%%%%%%%%%%%%%%%%%%%%%%%%%%%%%%%%%%
\subsection{Finite Variation}
\label{subsec:finite-variation-stochastic}
%%%%%%%%%%%%%%%%%%%%%%%%%%%%%%%%%%%%%%%%%%%%%%%%%%%%%%%%%%%%

A path \(A_t\) has finite variation on \([0,T]\) if

\[
\sup_{\Pi}
\sum_k
|A_{t_{k+1}}-A_{t_k}|
<
\infty,
\]

where the supremum is over all partitions.

Ordinary time integrals

\[
\int_0^t
\mu_s\,ds
\]

typically generate finite-variation paths.

Brownian motion does not have finite variation.

%%%%%%%%%%%%%%%%%%%%%%%%%%%%%%%%%%%%%%%%%%%%%%%%%%%%%%%%%%%%
\subsection{Semimartingale Differential Structure}
\label{subsec:semimartingale-differential}
%%%%%%%%%%%%%%%%%%%%%%%%%%%%%%%%%%%%%%%%%%%%%%%%%%%%%%%%%%%%

An It\^o process

\[
X_t
=
X_0
+
A_t
+
M_t
\]

contains:

\[
\boxed{
A_t
=
\int_0^t
\mu_s\,ds
}
\]

as a finite-variation component,

and

\[
\boxed{
M_t
=
\int_0^t
\sigma_s\,dB_s
}
\]

as a continuous local-martingale component.

This decomposition explains why stochastic differentials naturally split
into

\[
dt
\]

and

\[
dB_t
\]

terms.

%%%%%%%%%%%%%%%%%%%%%%%%%%%%%%%%%%%%%%%%%%%%%%%%%%%%%%%%%%%%
\subsection{Stochastic Calculus Workflow}
\label{subsec:stochastic-calculus-workflow}
%%%%%%%%%%%%%%%%%%%%%%%%%%%%%%%%%%%%%%%%%%%%%%%%%%%%%%%%%%%%

When solving an elementary stochastic-calculus problem:

\begin{enumerate}
    \item identify the stochastic process;
    \item identify the drift and diffusion terms;
    \item determine the filtration and adaptedness requirements;
    \item decide whether an It\^o integral is involved;
    \item use the It\^o isometry for second moments when appropriate;
    \item identify the function \(f(t,x)\) for It\^o's formula;
    \item compute \(f_t\), \(f_x\), and \(f_{xx}\);
    \item use
    \[
    (dB_t)^2=dt;
    \]
    \item separate the resulting drift and diffusion terms;
    \item check whether the stochastic integral has zero expectation;
    \item identify possible martingale structure;
    \item distinguish exact identities from heuristic differential
          notation.
\end{enumerate}

%%%%%%%%%%%%%%%%%%%%%%%%%%%%%%%%%%%%%%%%%%%%%%%%%%%%%%%%%%%%
\subsection{Common Errors}
\label{subsec:stochastic-calculus-common-errors}
%%%%%%%%%%%%%%%%%%%%%%%%%%%%%%%%%%%%%%%%%%%%%%%%%%%%%%%%%%%%

\subsubsection{Treating Brownian Motion as Differentiable}

The derivative

\[
\frac{dB_t}{dt}
\]

does not exist in the ordinary pathwise sense.

Stochastic differential notation should not be interpreted as ordinary
differentiation.

%%%%%%%%%%%%%%%%%%%%%%%%%%%%%%%%%%%%%%%%%%%%%%%%%%%%%%%%%%%%
\subsubsection{Using Ordinary Chain Rule Instead of It\^o's Formula}

For Brownian-driven processes, the ordinary chain rule misses the term

\[
\boxed{
\frac12
\sigma^2f_{xx}\,dt.
}
\]

%%%%%%%%%%%%%%%%%%%%%%%%%%%%%%%%%%%%%%%%%%%%%%%%%%%%%%%%%%%%
\subsubsection{Setting \((dB_t)^2=0\)}

Brownian quadratic variation gives

\[
\boxed{
(dB_t)^2=dt,
}
\]

in the stochastic differential bookkeeping sense.

%%%%%%%%%%%%%%%%%%%%%%%%%%%%%%%%%%%%%%%%%%%%%%%%%%%%%%%%%%%%
\subsubsection{Treating \((dB_t)^2=dt\) as Ordinary Algebra}

This notation summarizes quadratic variation.

It is not an ordinary equality between infinitesimals.

%%%%%%%%%%%%%%%%%%%%%%%%%%%%%%%%%%%%%%%%%%%%%%%%%%%%%%%%%%%%
\subsubsection{Forgetting Adaptedness}

The It\^o integral requires nonanticipating integrands under its standard
construction.

An integrand generally cannot depend on future Brownian information.

%%%%%%%%%%%%%%%%%%%%%%%%%%%%%%%%%%%%%%%%%%%%%%%%%%%%%%%%%%%%
\subsubsection{Assuming Every Stochastic Integral Is Gaussian}

If the integrand is deterministic, the Brownian stochastic integral is
Gaussian.

If the integrand itself is random, the resulting distribution need not
be Gaussian.

%%%%%%%%%%%%%%%%%%%%%%%%%%%%%%%%%%%%%%%%%%%%%%%%%%%%%%%%%%%%
\subsubsection{Assuming Every Local Martingale Is a Martingale}

Additional integrability conditions may be required.

%%%%%%%%%%%%%%%%%%%%%%%%%%%%%%%%%%%%%%%%%%%%%%%%%%%%%%%%%%%%
\subsubsection{Forgetting Quadratic Covariation in a Product Rule}

The correct stochastic product rule is

\[
d(XY)
=
X\,dY
+
Y\,dX
+
d[X,Y].
\]

%%%%%%%%%%%%%%%%%%%%%%%%%%%%%%%%%%%%%%%%%%%%%%%%%%%%%%%%%%%%
\subsubsection{Confusing It\^o and Stratonovich Integrals}

They use different limiting conventions and generally produce different
drift terms.

%%%%%%%%%%%%%%%%%%%%%%%%%%%%%%%%%%%%%%%%%%%%%%%%%%%%%%%%%%%%
\subsubsection{Applying Girsanov Without Integrability Conditions}

A stochastic exponential must satisfy appropriate conditions before it
can define a valid probability-measure change.

%%%%%%%%%%%%%%%%%%%%%%%%%%%%%%%%%%%%%%%%%%%%%%%%%%%%%%%%%%%%
\subsubsection{Assuming Zero Drift Automatically Means a True Martingale}

A driftless It\^o expression is often a local martingale.

True martingale status may require additional integrability.

%%%%%%%%%%%%%%%%%%%%%%%%%%%%%%%%%%%%%%%%%%%%%%%%%%%%%%%%%%%%
\subsection{Applications}
\label{subsec:stochastic-calculus-applications}
%%%%%%%%%%%%%%%%%%%%%%%%%%%%%%%%%%%%%%%%%%%%%%%%%%%%%%%%%%%%

\begin{applicationbox}

\textbf{Mathematical Finance.}
It\^o calculus is used for asset-price models, derivative pricing,
hedging, risk-neutral measures, and interest-rate models.

\medskip

\textbf{Physics.}
Brownian motion and stochastic differential equations model diffusion,
thermal noise, and random forcing.

\medskip

\textbf{Engineering.}
Stochastic calculus models noisy dynamical systems, filtering, control,
and signal propagation.

\medskip

\textbf{Biology.}
Random population growth, molecular motion, and diffusion-driven
biological processes can be modeled using stochastic calculus.

\medskip

\textbf{Stochastic Control.}
Dynamic optimization under uncertainty uses diffusion processes,
It\^o's formula, and Hamilton--Jacobi--Bellman equations.

\medskip

\textbf{Statistics.}
Continuous-time filtering and likelihood methods use martingales and
stochastic integrals.

\medskip

\textbf{Partial Differential Equations.}
Feynman--Kac formulas connect diffusion expectations with PDE solutions.

\medskip

\textbf{Probability Theory.}
Martingales, local time, semimartingales, and changes of measure are
central structures in modern continuous-time probability.

\end{applicationbox}

%%%%%%%%%%%%%%%%%%%%%%%%%%%%%%%%%%%%%%%%%%%%%%%%%%%%%%%%%%%%
\subsection{Exercises}
\label{subsec:stochastic-calculus-exercises}
%%%%%%%%%%%%%%%%%%%%%%%%%%%%%%%%%%%%%%%%%%%%%%%%%%%%%%%%%%%%

\begin{exercise}[1]
Explain why ordinary differentiation cannot be applied to Brownian
sample paths.
\end{exercise}

\begin{exercise}[1]
Define quadratic variation.
\end{exercise}

\begin{exercise}[1]
State the quadratic variation of standard Brownian motion.
\end{exercise}

\begin{exercise}[1]
State the symbolic It\^o rules for

\[
(dt)^2,
\qquad
dt\,dB_t,
\qquad
(dB_t)^2.
\]
\end{exercise}

\begin{exercise}[1]
Define a simple adapted process.
\end{exercise}

\begin{exercise}[1]
Define the It\^o integral of a simple adapted process.
\end{exercise}

\begin{exercise}[1]
State the It\^o isometry.
\end{exercise}

\begin{exercise}[1]
State It\^o's formula for

\[
dX_t
=
\mu_t\,dt+\sigma_t\,dB_t.
\]
\end{exercise}

\begin{exercise}[1]
Define quadratic covariation.
\end{exercise}

\begin{exercise}[1]
State the stochastic product rule.
\end{exercise}

\begin{exercise}[1]
What is a local martingale?
\end{exercise}

\begin{exercise}[1]
Distinguish It\^o and Stratonovich integration.
\end{exercise}

\begin{exercise}[2]
Compute the quadratic variation of a continuously differentiable
deterministic function.
\end{exercise}

\begin{exercise}[2]
For an equally spaced Brownian partition, compute the expected sum

\[
\mathbb{E}
\left[
\sum_k
(\Delta B_k)^2
\right].
\]
\end{exercise}

\begin{exercise}[2]
Let

\[
I
=
\int_0^T
1\,dB_t.
\]

Find \(I\).
\end{exercise}

\begin{exercise}[2]
Use the It\^o isometry to compute

\[
\mathbb{E}
\left[
\left(
\int_0^T
t\,dB_t
\right)^2
\right].
\]
\end{exercise}

\begin{exercise}[2]
Find the distribution of

\[
\int_0^T
2\,dB_t.
\]
\end{exercise}

\begin{exercise}[2]
Use It\^o's formula to compute

\[
d(B_t^2).
\]
\end{exercise}

\begin{exercise}[2]
Use It\^o's formula to compute

\[
d(B_t^3).
\]
\end{exercise}

\begin{exercise}[2]
Use It\^o's formula to compute

\[
d(e^{B_t}).
\]
\end{exercise}

\begin{exercise}[2]
Use It\^o's formula to compute

\[
d(\sin B_t).
\]
\end{exercise}

\begin{exercise}[2]
Use It\^o's formula to compute

\[
d(tB_t).
\]
\end{exercise}

\begin{exercise}[3]
Derive

\[
\int_0^t
B_s\,dB_s
=
\frac12
(B_t^2-t).
\]
\end{exercise}

\begin{exercise}[3]
Show that

\[
B_t^2-t
\]

is a martingale.
\end{exercise}

\begin{exercise}[3]
Show that

\[
B_t^3-3tB_t
\]

is a local martingale.
\end{exercise}

\begin{exercise}[3]
Show using It\^o's formula that

\[
\exp
\left(
\theta B_t-\frac12\theta^2t
\right)
\]

has zero drift.
\end{exercise}

\begin{exercise}[3]
Use the It\^o isometry to show that if

\[
M_t
=
\int_0^t
H_s\,dB_s,
\]

then

\[
\mathbb{E}[M_t^2]
=
\mathbb{E}
\left[
\int_0^t
H_s^2\,ds
\right].
\]
\end{exercise}

\begin{exercise}[3]
Show that

\[
[M]_t
=
\int_0^t
H_s^2\,ds
\]

for

\[
M_t
=
\int_0^t
H_s\,dB_s.
\]
\end{exercise}

\begin{exercise}[3]
Let \(B^{(1)}\) and \(B^{(2)}\) be independent Brownian motions. Show
that

\[
[B^{(1)},B^{(2)}]_t=0.
\]
\end{exercise}

\begin{exercise}[3]
Suppose Brownian motions have correlation \(\rho\). Use

\[
dB_t^{(1)}dB_t^{(2)}
=
\rho\,dt
\]

to derive the stochastic differential of

\[
B_t^{(1)}B_t^{(2)}.
\]
\end{exercise}

\begin{exercise}[3]
Use stochastic integration by parts to derive

\[
\int_0^t
s\,dB_s
=
tB_t
-
\int_0^t
B_s\,ds.
\]
\end{exercise}

\begin{exercise}[3]
Suppose

\[
dX_t
=
a\,dt+b\,dB_t.
\]

Use It\^o's formula to find

\[
d(X_t^2).
\]
\end{exercise}

\begin{exercise}[3]
For the same process, find

\[
d(e^{X_t}).
\]
\end{exercise}

\begin{exercise}[3]
Suppose

\[
dX_t
=
\mu X_t\,dt
+
\sigma X_t\,dB_t.
\]

Derive

\[
d(\ln X_t).
\]
\end{exercise}

\begin{exercise}[3]
Show that the quadratic variation of

\[
X_t
=
\int_0^t
\sigma_s\,dB_s
\]

is

\[
\int_0^t
\sigma_s^2\,ds.
\]
\end{exercise}

\begin{exercise}[3]
Use It\^o's formula to derive the generator of

\[
dX_t
=
\mu(X_t)\,dt
+
\sigma(X_t)\,dB_t.
\]
\end{exercise}

\begin{exercise}[3]
Show that if

\[
f_t+\frac12f_{xx}=0,
\]

then \(f(t,B_t)\) has zero It\^o drift.
\end{exercise}

\begin{exercise}[3]
Compare

\[
\int_0^tB_s\,dB_s
\]

and

\[
\int_0^tB_s\circ dB_s.
\]
\end{exercise}

\begin{exercise}[3]
Convert

\[
dX_t
=
a(X_t)\,dt
+
b(X_t)\circ dB_t
\]

from Stratonovich to It\^o form.
\end{exercise}

\begin{exercise}[4]
Construct the It\^o integral rigorously using \(L^2\) completion.
\end{exercise}

\begin{exercise}[4]
Prove the It\^o isometry first for simple processes and then extend it by
density.
\end{exercise}

\begin{exercise}[4]
Study continuous local martingales and prove that their quadratic
variation is increasing.
\end{exercise}

\begin{exercise}[4]
Investigate semimartingales and the general stochastic integration
framework.
\end{exercise}

\begin{exercise}[4]
Study multidimensional It\^o's formula.
\end{exercise}

\begin{exercise}[4]
Investigate stochastic exponentials and Novikov's condition.
\end{exercise}

\begin{exercise}[4]
Study Girsanov's theorem and prove a one-dimensional Brownian drift-change
version.
\end{exercise}

\begin{exercise}[4]
Investigate the Feynman--Kac formula.
\end{exercise}

\begin{exercise}[4]
Study local time and Tanaka's formula.
\end{exercise}

\begin{exercise}[4]
Investigate the occupation-time formula for Brownian motion.
\end{exercise}

\begin{exercise}[4]
Study the relationship between It\^o and Stratonovich integrals in
multiple dimensions.
\end{exercise}

\begin{exercise}[4]
Investigate stochastic calculus on manifolds.
\end{exercise}

\begin{exercise}[4]
Study martingale representation for Brownian filtrations.
\end{exercise}

\begin{exercise}[4]
Investigate stochastic integration with respect to general
semimartingales.
\end{exercise}

%%%%%%%%%%%%%%%%%%%%%%%%%%%%%%%%%%%%%%%%%%%%%%%%%%%%%%%%%%%%
\subsection{Conceptual Summary}
\label{subsec:stochastic-calculus-summary}
%%%%%%%%%%%%%%%%%%%%%%%%%%%%%%%%%%%%%%%%%%%%%%%%%%%%%%%%%%%%

Brownian motion is continuous but almost surely nowhere differentiable.

Its quadratic variation is

\[
\boxed{
[B]_t=t.
}
\]

This motivates the stochastic differential rule

\[
\boxed{
(dB_t)^2=dt.
}
\]

For a simple adapted process,

\[
\boxed{
\int_0^T
H_t\,dB_t
=
\sum_k
H_k
(B_{t_{k+1}}-B_{t_k}).
}
\]

The It\^o isometry is

\[
\boxed{
\mathbb{E}
\left[
\left(
\int_0^T
H_t\,dB_t
\right)^2
\right]
=
\mathbb{E}
\left[
\int_0^T
H_t^2\,dt
\right].
}
\]

For

\[
dX_t
=
\mu_t\,dt
+
\sigma_t\,dB_t,
\]

It\^o's formula gives

\[
\boxed{
df(t,X_t)
=
\left(
f_t
+
\mu_tf_x
+
\frac12\sigma_t^2f_{xx}
\right)dt
+
\sigma_tf_x\,dB_t.
}
\]

The stochastic product rule is

\[
\boxed{
d(X_tY_t)
=
X_t\,dY_t
+
Y_t\,dX_t
+
d[X,Y]_t.
}
\]

For an It\^o process,

\[
\boxed{
d[X]_t
=
\sigma_t^2\,dt.
}
\]

The exponential process

\[
\boxed{
\exp
\left(
\theta B_t
-
\frac12\theta^2t
\right)
}
\]

is a fundamental martingale.

It\^o and Stratonovich integration differ by a quadratic-covariation
correction:

\[
\boxed{
\int H\circ dB
=
\int H\,dB
+
\frac12[H,B].
}
\]

Stochastic calculus therefore replaces ordinary differential calculus
with a framework that accounts for

\[
\boxed{
\text{random variation}
+
\text{quadratic variation}
+
\text{martingales}
+
\text{nondifferentiable paths}.
}
\]

%%%%%%%%%%%%%%%%%%%%%%%%%%%%%%%%%%%%%%%%%%%%%%%%%%%%%%%%%%%%
\subsection{Section Connections}
\label{subsec:stochastic-calculus-connections}
%%%%%%%%%%%%%%%%%%%%%%%%%%%%%%%%%%%%%%%%%%%%%%%%%%%%%%%%%%%%

\chapterconnection{
Stochastic Calculus develops the continuous-time integration and chain
rule machinery required for Brownian-driven models introduced in
Section~\ref{sec:stochastic-processes}. Measure-Theoretic Probability
provides the filtrations, conditional expectations, \(L^2\) spaces, and
martingale concepts used to construct the It\^o integral rigorously.
Quadratic variation explains why It\^o's formula differs from ordinary
calculus, while generators and Feynman--Kac ideas connect stochastic
processes with partial differential equations. The next section,
Stochastic Differential Equations, uses It\^o integrals and It\^o's
formula to define, solve, analyze, and simulate differential equations
driven by Brownian motion.
}

%%%%%%%%%%%%%%%%%%%%%%%%%%%%%%%%%%%%%%%%%%%%%%%%%%%%%%%%%%%%
% SECTION SOURCE TRACKING
%%%%%%%%%%%%%%%%%%%%%%%%%%%%%%%%%%%%%%%%%%%%%%%%%%%%%%%%%%%%

%------------------------------------------------------------
% 7.13 Stochastic Calculus
%------------------------------------------------------------
%
% Source basis:
%   - Standard Brownian-motion theory.
%   - Standard Itô stochastic-integration theory.
%   - Standard martingale and semimartingale calculus.
%
% Used for:
%   - Brownian path irregularity;
%   - Brownian increment scaling;
%   - quadratic variation;
%   - Brownian quadratic variation;
%   - stochastic differential rules;
%   - simple adapted processes;
%   - Itô integrals;
%   - L2 construction of stochastic integrals;
%   - Itô isometry;
%   - mean and variance of stochastic integrals;
%   - stochastic-integral martingales;
%   - quadratic covariation;
%   - correlated Brownian motions;
%   - Itô's formula;
%   - Brownian-square example;
%   - exponential martingales;
%   - time-dependent Itô formula;
%   - heat-equation connection;
%   - diffusion generators;
%   - Dynkin's formula preview;
%   - stochastic product rule;
%   - stochastic integration by parts;
%   - multidimensional Itô formula;
%   - multidimensional Brownian motion;
%   - Itô processes;
%   - drift and diffusion;
%   - quadratic variation of Itô processes;
%   - local martingales;
%   - stochastic exponentials;
%   - Novikov condition preview;
%   - Stratonovich integration;
%   - Itô--Stratonovich conversion;
%   - Girsanov theorem preview;
%   - probability-measure change;
%   - Feynman--Kac preview;
%   - local time and Tanaka formula preview;
%   - semimartingale preview;
%   - finite variation;
%   - applications and exercises.
%
% Notes:
%   - Stochastic Differential Equations are developed next in
%     Section 7.14.
%   - Full SDE existence, uniqueness, explicit solutions, numerical
%     schemes, and stability analysis are deferred to Section 7.14.
%   - Brownian motion was introduced in Section 7.12 and is used here as
%     the canonical stochastic integrator.
%   - Measure-theoretic foundations appear in Section 7.11.
%
%%%%%%%%%%%%%%%%%%%%%%%%%%%%%%%%%%%%%%%%%%%%%%%%%%%%%%%%%%%%